% !TeX program = xelatex
% Компилировать XeLaTeX (в Overleaf: Menu -> Compiler -> XeLaTeX)
\documentclass[14pt,a4paper]{extarticle}


\usepackage[left=2cm,right=2cm,top=2cm,bottom=2cm]{geometry}
\usepackage{fontspec}
\defaultfontfeatures{Ligatures=TeX}
\usepackage{polyglossia}
\setdefaultlanguage{russian}
\setotherlanguage{english}
\IfFontExistsTF{Times New Roman}
  {\setmainfont{Times New Roman}\newfontfamily\cyrillicfont{Times New Roman}}
  {\setmainfont{Liberation Serif}\newfontfamily\cyrillicfont{Liberation Serif}}
\IfFontExistsTF{DejaVu Sans Mono}
  {\setmonofont[Scale=0.85]{DejaVu Sans Mono}\newfontfamily\cyrillicfonttt[Scale=0.85]{DejaVu Sans Mono}}
  {}

\usepackage{amsmath,amssymb}
\usepackage{graphicx}
\usepackage{float}
\usepackage{array,tabularx,booktabs,multirow}
\usepackage{caption}
\usepackage{enumitem}
\usepackage{indentfirst}
\usepackage{setspace}
\usepackage{titlesec}
\usepackage{tocloft}
\usepackage{listings}
\usepackage{xcolor}
\usepackage{tikz}
\usetikzlibrary{positioning,shapes.geometric,arrows.meta,calc}
\usepackage[hidelinks]{hyperref}

\onehalfspacing
\setlength{\parindent}{1.25cm}
\graphicspath{{images/}}

\captionsetup{labelsep=endash,font=small}
\captionsetup[table]{justification=raggedright,singlelinecheck=false}
\AtBeginDocument{%
  \renewcommand{\figurename}{Рисунок}%
  \renewcommand{\tablename}{Таблица}%
  
}

\titleformat{\section}{\normalsize\bfseries\centering}{\thesection}{0.6em}{}
\titleformat{\subsection}{\normalsize\bfseries}{\thesubsection}{0.6em}{}
\titlespacing*{\section}{0pt}{0pt}{1.5em}
\titlespacing*{\subsection}{\parindent}{1.2em}{0.8em}
\renewcommand{\cfttoctitlefont}{\hfill\normalsize\bfseries}
\renewcommand{\cftaftertoctitle}{\hfill\null}
\gappto\captionsrussian{\renewcommand{\contentsname}{СОДЕРЖАНИЕ}\renewcommand{\figurename}{Рисунок}\renewcommand{\tablename}{Таблица}}
\renewcommand{\cftsecleader}{\cftdotfill{\cftdotsep}}
\renewcommand{\cftsecfont}{\normalfont}
\renewcommand{\cftsecpagefont}{\normalfont}

\setlist{nosep,leftmargin=*,labelindent=\parindent}

\lstset{
  basicstyle=\ttfamily\scriptsize,
  columns=fullflexible,
  breaklines=true,
  frame=single,
  numbers=left,
  numberstyle=\tiny,
  keywordstyle=\color{blue!70!black}\bfseries,
  commentstyle=\color{green!45!black},
  stringstyle=\color{red!60!black},
  extendedchars=true,
  language=Python
}

\tikzset{
  >={Stealth[length=2.2mm]},
  block/.style={draw,rectangle,minimum height=9mm,minimum width=11mm,align=center},
  sum/.style={draw,circle,minimum size=5mm,inner sep=0pt},
  gainL/.style={draw,isosceles triangle,shape border rotate=180,isosceles triangle apex angle=60,
                minimum height=8mm,inner sep=1pt,font=\small},
  gainR/.style={draw,isosceles triangle,isosceles triangle apex angle=60,
                minimum height=8mm,inner sep=1pt,font=\small},
  dot/.style={circle,fill,inner sep=1.2pt},
}

\begin{document}

% =================== ТИТУЛЬНЫЙ ЛИСТ ===================
\begin{titlepage}
\begin{center}
\small
МИНИСТЕРСТВО НАУКИ И ВЫСШЕГО ОБРАЗОВАНИЯ РОССИЙСКОЙ\\
ФЕДЕРАЦИИ\\[0.3em]
ФЕДЕРАЛЬНОЕ ГОСУДАРСТВЕННОЕ АВТОНОМНОЕ ОБРАЗОВАТЕЛЬНОЕ\\
УЧРЕЖДЕНИЕ ВЫСШЕГО ОБРАЗОВАНИЯ\\[0.3em]
«Национальный исследовательский университет ИТМО»\\[0.8em]
\normalsize
ФАКУЛЬТЕТ СИСТЕМ УПРАВЛЕНИЯ И РОБОТОТЕХНИКИ

\vfill

{\bfseries ОТЧЕТ ПО ЛАБОРАТОРНОЙ РАБОТЕ №2}\\[0.5em]
по дисциплине «Линейные системы автоматического управления»\\[0.5em]
по теме:\\
{\bfseries Свободное движение и устойчивость}

\vfill

\begin{flushright}
\begin{tabular}{l}
Выполнили:\\
Басиони Абдельсалам Али Абдельсалам\\
Боврджи Тамим Махер Мохаммед\\
Май Хоанг Куан\\[1em]
Преподаватель:\\
Пашенко Артем Витальевич
\end{tabular}
\end{flushright}

\vfill

Санкт-Петербург\\
2026
\end{center}
\end{titlepage}

\setcounter{page}{2}

\newpage

\tableofcontents

\newpage
% =================== ЗАДАНИЕ 1 ===================
\section{СВОБОДНОЕ ДВИЖЕНИЕ}

Рассмотрим линейную систему второго порядка, описываемую дифференциальным уравнением
\begin{equation}\label{eq:sys2}
  \ddot y + a_1\dot y + a_0 y = u
\end{equation}
с начальными условиями $y(0)=y_0$, $\dot y(0)=\dot y_0$.
Так как исследуется свободное движение, внешнее воздействие принимается равным нулю: $u(t)\equiv 0$.

\subsection{Схема моделирования}

Выразим старшую производную: $\ddot y = u - a_1\dot y - a_0 y$.
Схема моделирования, построенная из блоков элементарных операций (сумматор, усилители, интеграторы), приведена на рисунке~\ref{fig:scheme1}.
Начальные условия $\dot y(0)$ и $y(0)$ задаются в первом и втором интеграторах соответственно.


\begin{figure}[H]
    \centering
    \includegraphics[width=1.1\linewidth]{scheme1.png}
    \caption{Схема моделирования системы второго порядка}
    \label{fig:scheme1}
\end{figure}


\subsection{Аналитические расчёты}

Характеристическое уравнение системы \eqref{eq:sys2}:
\[
  \lambda^2 + a_1\lambda + a_0 = 0 .
\]
По теореме Виета коэффициенты выражаются через корни:
\[
  a_1 = -(\lambda_1+\lambda_2), \qquad a_0 = \lambda_1\lambda_2 .
\]
Свободная составляющая движения является решением однородного уравнения и в зависимости от вида корней имеет вид:
\begin{itemize}
  \item различные вещественные корни:
    \[ y_{\text{св}}(t) = c_1e^{\lambda_1 t} + c_2e^{\lambda_2 t}; \]
  \item комплексно-сопряжённые корни $\lambda_{1,2}=\alpha\pm j\beta$:
    \[ y_{\text{св}}(t) = e^{\alpha t}\left(c_1\cos\beta t + c_2\sin\beta t\right). \]
\end{itemize}
Константы $c_1$, $c_2$ находятся из начальных условий $y(0)=y_0$, $\dot y(0)=\dot y_0$.


\paragraph{Эксперимент 1.} $\lambda_1=-3$, $\lambda_2=-2$, $y_0=1$, $\dot y_0=0$.

Коэффициенты:
\[
  a_1 = -(\lambda_1 + \lambda_2) = -((-3) + (-2)) = 5
\]
\[
  a_0 = \lambda_1 \cdot \lambda_2 = (-3) \cdot (-2) = 6
\]

Дифференциальное уравнение:
\begin{equation}
  \ddot y + 5\dot y + 6y = 0
\end{equation}

Общее решение для различных вещественных корней:
\[
  y(t) = C_1 e^{-3t} + C_2 e^{-2t}
\]

Производная:
\[
  \dot y(t) = -3C_1 e^{-3t} - 2C_2 e^{-2t}
\]

Начальные условия при $t = 0$:
\[
  \begin{cases}
    y(0) = C_1 + C_2 = 1 \\
    \dot y(0) = -3C_1 - 2C_2 = 0
  \end{cases}
  \quad\Rightarrow\quad
  \begin{cases}
    C_1 = -2 \\
    C_2 = 3
  \end{cases}
\]

Итоговое решение:
\begin{equation}
  y_{\text{св}}(t) = 3e^{-2t} - 2e^{-3t}
\end{equation}


\paragraph{Эксперимент 2.} $\lambda_{1,2}=-0.5\pm j3$, $y_0=1$, $\dot y_0=0$.

Коэффициенты:
\[
  a_1 = -(\lambda_1 + \lambda_2) = -((-0.5 + j3) + (-0.5 - j3)) = 1
\]
\[
  a_0 = \lambda_1 \cdot \lambda_2 = (-0.5 + j3)(-0.5 - j3) = (-0.5)^2 + 3^2 = 9.25
\]

Дифференциальное уравнение:
\begin{equation}
  \ddot y + \dot y + 9.25y = 0
\end{equation}

Общее решение для комплексно-сопряженных корней:
\[
  y(t) = e^{-0.5t}(C_1 \cos 3t + C_2 \sin 3t)
\]

Производная:
\[
  \dot y(t) = e^{-0.5t}\left[(C_1 \cos 3t + C_2 \sin 3t)(-0.5) + (-3C_1 \sin 3t + 3C_2 \cos 3t)\right]
\]

Начальные условия при $t = 0$:
\[
  \begin{cases}
    y(0) = C_1 = 1 \\
    \dot y(0) = -0.5 \cdot C_1 + 3 \cdot C_2 = 0
  \end{cases}
  \quad\Rightarrow\quad
  \begin{cases}
    C_1 = 1 \\
    C_2 = \frac{1}{6}
  \end{cases}
\]

Итоговое решение:
\begin{equation}
  y_{\text{св}}(t) = e^{-0.5t}\left(\cos 3t + \frac{1}{6}\sin 3t\right)
\end{equation}


\paragraph{Эксперимент 3.} $\lambda_{1,2}=\pm j3$, $y_0=1$, $\dot y_0=0$.

Коэффициенты:
\[
  a_1 = -(j3 + (-j3)) = 0
\]
\[
  a_0 = (j3)(-j3) = 3^2 = 9
\]

Дифференциальное уравнение:
\begin{equation}
  \ddot y + 9y = 0
\end{equation}

Общее решение для чисто мнимых корней:
\[
  y(t) = C_1 \cos 3t + C_2 \sin 3t
\]

Производная:
\[
  \dot y(t) = -3C_1 \sin 3t + 3C_2 \cos 3t
\]

Начальные условия при $t = 0$:
\[
  \begin{cases}
    y(0) = C_1 = 1 \\
    \dot y(0) = 3C_2 = 0
  \end{cases}
  \quad\Rightarrow\quad
  \begin{cases}
    C_1 = 1 \\
    C_2 = 0
  \end{cases}
\]

Итоговое решение:
\begin{equation}
  y_{\text{св}}(t) = \cos 3t
\end{equation}


\paragraph{Эксперимент 4.} $\lambda_{1,2}=0.5\pm j3$, $y_0=0.05$, $\dot y_0=0$.

Коэффициенты:
\[
  a_1 = -((0.5 + j3) + (0.5 - j3)) = -1
\]
\[
  a_0 = (0.5 + j3)(0.5 - j3) = 0.5^2 + 3^2 = 9.25
\]

Дифференциальное уравнение:
\begin{equation}
  \ddot y - \dot y + 9.25y = 0
\end{equation}

Общее решение для комплексно-сопряженных корней с положительной вещественной частью:
\[
  y(t) = e^{0.5t}(C_1 \cos 3t + C_2 \sin 3t)
\]

Производная:
\[
  \dot y(t) = e^{0.5t}\left[(C_1 \cos 3t + C_2 \sin 3t)(0.5) + (-3C_1 \sin 3t + 3C_2 \cos 3t)\right]
\]

Начальные условия при $t = 0$:
\[
  \begin{cases}
    y(0) = C_1 = 0.05 \\
    \dot y(0) = 0.5 \cdot C_1 + 3 \cdot C_2 = 0
  \end{cases}
  \quad\Rightarrow\quad
  \begin{cases}
    C_1 = 0.05 \\
    C_2 = -\frac{1}{120}
  \end{cases}
\]

Итоговое решение:
\begin{equation}
  y_{\text{св}}(t) = e^{0.5t}\left(0.05 \cos 3t - \frac{1}{120} \sin 3t\right)
\end{equation}


\paragraph{Эксперимент 5.} $\lambda_1=3$, $\lambda_2=2$, $y_0=0.05$, $\dot y_0=0$.

Коэффициенты:
\[
  a_1 = -(3 + 2) = -5
\]
\[
  a_0 = 3 \cdot 2 = 6
\]

Дифференциальное уравнение:
\begin{equation}
  \ddot y - 5\dot y + 6y = 0
\end{equation}

Общее решение для различных вещественных корней:
\[
  y(t) = C_1 e^{3t} + C_2 e^{2t}
\]

Производная:
\[
  \dot y(t) = 3C_1 e^{3t} + 2C_2 e^{2t}
\]

Начальные условия при $t = 0$:
\[
  \begin{cases}
    y(0) = C_1 + C_2 = 0.05 \\
    \dot y(0) = 3C_1 + 2C_2 = 0
  \end{cases}
  \quad\Rightarrow\quad
  \begin{cases}
    C_1 = -0.1 \\
    C_2 = 0.15
  \end{cases}
\]

Итоговое решение:
\begin{equation}
  y_{\text{св}}(t) = 0.15e^{2t} - 0.1e^{3t}
\end{equation}


\paragraph{Эксперимент 6.} $\lambda_1=-0.9$, $\lambda_2=0.9$, $y_0=0$, $\dot y_0=0.1$.

Коэффициенты:
\[
  a_1 = -((-0.9) + 0.9) = 0
\]
\[
  a_0 = (-0.9) \cdot 0.9 = -0.81
\]

Дифференциальное уравнение:
\begin{equation}
  \ddot y - 0.81y = 0
\end{equation}

Общее решение для вещественных корней:
\[
  y(t) = C_1 e^{-0.9t} + C_2 e^{0.9t}
\]

Производная:
\[
  \dot y(t) = -0.9C_1 e^{-0.9t} + 0.9C_2 e^{0.9t}
\]

Начальные условия при $t = 0$:
\[
  \begin{cases}
    y(0) = C_1 + C_2 = 0 \\
    \dot y(0) = -0.9C_1 + 0.9C_2 = 0.1
  \end{cases}
  \quad\Rightarrow\quad
  \begin{cases}
    C_1 = -\frac{1}{18} \\
    C_2 = \frac{1}{18}
  \end{cases}
\]

Итоговое решение:
\begin{equation}
  y_{\text{св}}(t) = \frac{1}{18}\left(e^{0.9t} - e^{-0.9t}\right) = \frac{1}{9}\operatorname{sh}(0.9t)
\end{equation}



\subsection{Результаты вычислений}

Результаты расчётов сведены в таблицу~\ref{tab:task1}. Тип устойчивости определён по корневому критерию.

\begin{table}[H]
\caption{Результаты вычислений для Задания 1}\label{tab:task1}
\centering\footnotesize
\renewcommand{\arraystretch}{1.5}
\setlength{\tabcolsep}{3pt}
\begin{tabularx}{\textwidth}{|c|c|c|c|c|c|c|>{\centering\arraybackslash}X|>{\centering\arraybackslash}p{2.6cm}|}
\hline
№ & $\lambda_1$ & $\lambda_2$ & $a_0$ & $a_1$ & $y(0)$ & $\dot y(0)$ & $y_{\text{св}}(t)$ & Тип устойчивости \\ \hline
1 & $-3$ & $-2$ & $6$ & $5$ & $1$ & $0$ & $3e^{-2t}-2e^{-3t}$ & асимптотически устойчива \\ \hline
2 & $-0.5+j3$ & $-0.5-j3$ & $9.25$ & $1$ & $1$ & $0$ & $e^{-0.5t}\!\left(\cos 3t+\frac16\sin 3t\right)$ & асимптотически устойчива \\ \hline
3 & $j3$ & $-j3$ & $9$ & $0$ & $1$ & $0$ & $\cos 3t$ & на границе устойчивости \\ \hline
4 & $0.5+j3$ & $0.5-j3$ & $9.25$ & $-1$ & $0.05$ & $0$ & $e^{0.5t}\!\left(0.05\cos 3t-\frac{1}{120}\sin 3t\right)$ & неустойчива \\ \hline
5 & $3$ & $2$ & $6$ & $-5$ & $0.05$ & $0$ & $0.15e^{2t}-0.1e^{3t}$ & неустойчива \\ \hline
6 & $-0.9$ & $0.9$ & $-0.81$ & $0$ & $0$ & $0.1$ & $\frac{1}{18}\left(e^{0.9t}-e^{-0.9t}\right)$ & неустойчива \\ \hline
\end{tabularx}
\end{table}


\newpage
\subsection{Результаты моделирования}

На рисунках~\ref{fig:t1e1}--\ref{fig:t1e6} приведены сравнительные графики выхода модели $y(t)$ и аналитически рассчитанной свободной составляющей $y_{\text{св}}(t)$. На всех графиках сплошной красной линией показан выход модели $y(t)$, синей пунктирной линией~--- аналитическое решение $y_{\text{св}}(t)$.

\newcommand{\taskonefig}[3]{%
\begin{figure}[H]
\centering
\includegraphics[width=0.92\textwidth]{#1}
\caption{#2}\label{#3}
\end{figure}}

\taskonefig{experiment_1}{Эксперимент 1: $\lambda_{1,2}=-3,\,-2$}{fig:t1e1}
\taskonefig{experiment_2}{Эксперимент 2: $\lambda_{1,2}=-0.5\pm j3$}{fig:t1e2}
\taskonefig{experiment_3}{Эксперимент 3: $\lambda_{1,2}=\pm j3$}{fig:t1e3}
\taskonefig{experiment_4}{Эксперимент 4: $\lambda_{1,2}=0.5\pm j3$}{fig:t1e4}
\taskonefig{experiment_5}{Эксперимент 5: $\lambda_{1,2}=3,\,2$ (нулевые начальные условия)}{fig:t1e5}
\taskonefig{experiment_6}{Эксперимент 6: $\lambda_{1,2}=-0.9,\,0.9$}{fig:t1e6}


\newpage
\subsection{Выводы по заданию 1}

Было исследовано свободное движение линейных систем второго порядка. Установлена прямая связь между расположением корней характеристического уравнения и устойчивостью системы. Результаты аналитического расчета полностью подтвердились данными компьютерного моделирования.

\begin{enumerate}
    \item \textbf{Эксперимент 1:} Система асимптотически устойчива ($\lambda_{1,2}$ — вещественные отрицательные). На графике видно, что свободное движение плавно затухает и стремится к нулю ($\operatorname{Re}\lambda_1 < 0, \operatorname{Re}\lambda_2 < 0$).
    
    \item \textbf{Эксперимент 2:} Система асимптотически устойчива ($\lambda_{1,2}$ — комплексно-сопряженные с отрицательной вещественной частью). На графике наблюдаются затухающие колебания, амплитуда которых со временем стремится к нулю ($\operatorname{Re}\lambda_1 < 0, \operatorname{Re}\lambda_2 < 0$).
    
    \item \textbf{Эксперимент 3:} Система находится на границе устойчивости ($\lambda_{1,2}$ — чисто мнимые корни). На графике возникают незатухающие гармонические колебания с постоянной амплитудой ($\operatorname{Re}\lambda_1 = \operatorname{Re}\lambda_2 = 0, \operatorname{Im}\lambda_{1,2} \neq 0$).
    
    \item \textbf{Эксперимент 4:} Система неустойчива ($\lambda_{1,2}$ — комплексно-сопряженные с положительной вещественной частью). На графике наблюдаются расходящиеся колебания с увеличивающейся амплитудой ($\operatorname{Re}\lambda_1 > 0, \operatorname{Re}\lambda_2 > 0$).
    
    \item \textbf{Эксперимент 5:} Система неустойчива ($\lambda_{1,2}$ — положительные вещественные корни). Свободное движение экспоненциально возрастает по модулю ($\operatorname{Re}\lambda_1 > 0, \operatorname{Re}\lambda_2 > 0$).
    
    \item \textbf{Эксперимент 6:} Система неустойчива (корни имеют разные знаки, один из них положительный). Свободное движение неограниченно возрастает со временем ($\operatorname{Re}\lambda_2 > 0$).
\end{enumerate}


\newpage
\section{ОБЛАСТЬ УСТОЙЧИВОСТИ}

\subsection{Схема моделирования}

В соответствии с заданием собрана схема моделирования линейной системы третьего порядка (см. рисунок~\ref{fig:scheme2}), включающая в себя пропорциональный регулятор с коэффициентом $K$, два апериодических звена с постоянными времени $T_1$ и $T_2$, а также интегрирующее звено.

\begin{figure}[H]
    \centering
    \includegraphics[width=1.1\linewidth]{images/scheme2.png} 
    \caption{Структурная схема линейной системы третьего порядка}
    \label{fig:scheme2}
\end{figure}

Значения постоянных времени $T_1$ и $T_2$ заданы на основе первого набора корней характеристического уравнения из Задания 1 ($\lambda_1 = -3$, $\lambda_2 = -2$). Полюсы передаточных функций звеньев соответствуют данным корням:
\[
  T_1 = -\frac{1}{\lambda_1} = \frac{1}{3} \approx 0.333\text{ с}, \qquad
  T_2 = -\frac{1}{\lambda_2} = \frac{1}{2} = 0.5\text{ с}.
\]

\subsection{Аналитическое определение области устойчивости}

Для построенной структурной схемы (см. рисунок~\ref{fig:scheme2}) передаточная функция разомкнутой системы имеет вид:
\[
  W(s) = \frac{K}{s(T_1s + 1)(T_2s + 1)}.
\]
Характеристическое уравнение замкнутой системы определяется как $1 + W(s) = 0$:
\[
  T_1 T_2 s^3 + (T_1 + T_2) s^2 + s + K = 0.
\]
Коэффициенты характеристического полинома:
\[
  a_3 = T_1 T_2, \quad a_2 = T_1 + T_2, \quad a_1 = 1, \quad a_0 = K.
\]

Для исследования устойчивости применим критерий Гурвица. Главные угловые миноры матрицы Гурвица для системы 3-го порядка:
\[
  \Delta_1 = a_2 = T_1 + T_2 > 0,
\]
\[
  \Delta_2 = a_2 a_1 - a_3 a_0 = (T_1 + T_2) - (T_1 T_2)K > 0.
\]

При фиксированном значении $T_2 = 0.5$ с, граница устойчивости ($\Delta_2 = 0$) задается уравнением:
\[
  K = \frac{1}{T_1} + \frac{1}{T_2} = \frac{1}{T_1} + 2.
\]

Таким образом, область устойчивости в пространстве параметров $(K, T_1)$ определяется неравенствами:
\[
  T_1 > 0, \quad 0 < K < \frac{1}{T_1} + 2.
\]

На рисунке ниже представлено графическое изображение границы и области устойчивости системы на плоскости параметров $(K, T_1)$.

\begin{figure}[H]
    \centering
    \includegraphics[width=0.8\linewidth]{images/stability_region.png}
    \caption{Область устойчивости в пространстве параметров $K$ и $T_1$}
    \label{fig:stability_region}
\end{figure}


\subsection{Выбор параметров для моделирования переходных процессов}

Для исследования динамики системы зададим три набора параметров ($T_1 = 1/3$ с, $T_2 = 0.5$ с):
\begin{table}[H]
    \centering
    \caption{Наборы параметров для моделирования}
    \renewcommand{\arraystretch}{1.4} 
    \setlength{\tabcolsep}{10pt}     
    \begin{tabular}{|p{7cm}|c|c|c|}
        \hline
        \textbf{Состояние системы} & \textbf{$K$} & \textbf{$T_1$, с} & \textbf{$T_2$, с} \\ \hline
        а) Асимптотически устойчивая & $2$ & $1/3$ & $0.5$ \\ \hline
        б) На границе устойчивости   & $5$ & $1/3$ & $0.5$ \\ \hline
        в) Неустойчивая система      & $7$ & $1/3$ & $0.5$ \\ \hline
    \end{tabular}
    \label{tab:params}
\end{table}

\begin{enumerate}
    \item[а)] \textbf{Устойчивая система ($K = 2$):} условие $2 < 5$ выполнено, точка внутри области.
    \item[б)] \textbf{На границе ($K = 5$):} критическое значение ($5 = 5$), точка лежит на границе.
    \item[в)] \textbf{Неустойчивая система ($K = 7$):} условие не выполнено ($7 > 5$), точка в зоне неустойчивости.
\end{enumerate}


\subsection{Моделирование свободного движения системы}

Для каждого из трех выбранных наборов параметров ($T_1 = 1/3$ с, $T_2 = 0.5$ с) было выполнено моделирование свободного движения системы при нулевом внешнем воздействии ($g(t) = 0$) и ненулевых начальных условиях: $y(0) = 1$, $\dot{y}(0) = 0$, $\ddot{y}(0) = 0$.

\begin{enumerate}
    \item[а)] \textbf{Асимптотически устойчивая система ($K = 2$):}
    На графике (рисунок~\ref{fig:free_stable}) видно, что благодаря отрицательной обратной связи и нахождению внутри области устойчивости, свободные колебания затухают со временем, и система возвращается в нулевое положение равновесия.
    \begin{figure}[H]
        \centering
        \includegraphics[width=0.85\linewidth]{free_stable.png}
        \caption{Свободное движение асимптотически устойчивой системы ($K = 2, y(0) = 1$)}
        \label{fig:free_stable}
    \end{figure}

    \item[б)] \textbf{Система на границе устойчивости ($K = 5$):}
    На графике (рисунок~\ref{fig:free_boundary}) наблюдаются незатухающие гармонические колебания постоянной амплитуды, так как корни характеристического уравнения лежат на мнимой оси.
    \begin{figure}[H]
        \centering
        \includegraphics[width=0.85\linewidth]{free_boundary.png}
        \caption{Свободное движение на границе устойчивости ($K = 5, y(0) = 1$)}
        \label{fig:free_boundary}
    \end{figure}

    \item[в)] \textbf{Неустойчивая система ($K = 7$):}
    На графике (рисунок~\ref{fig:free_unstable}) зафиксирован расходящийся колебательный процесс с экспоненциально нарастающей амплитудой, что подтверждает неустойчивость системы при данном коэффициенте передачи.
    \begin{figure}[H]
        \centering
        \includegraphics[width=0.85\linewidth]{free_unstable.png}
        \caption{Свободное движение неустойчивой системы ($K = 7, y(0) = 1$)}
        \label{fig:free_unstable}
    \end{figure}
\end{enumerate}

\subsection{Выводы по заданию 2}

На основе аналитического расчета и компьютерного моделирования линейной системы третьего порядка можно сделать следующие выводы:
\begin{enumerate}
    \item Анализ по критерию Гурвица позволил определить критическое значение коэффициента передачи регулятора $K_{кр} = 5$ для заданных параметров $T_1 = 1/3$ с и $T_2 = 0.5$ с.
    \item Моделирование свободного движения при начальных условиях $y(0) = 1$, $\dot{y}(0) = 0$, $\ddot{y}(0) = 0$ полностью подтвердило теоретические расчёты:
    \begin{itemize}
        \item При $K = 2$ система асимптотически устойчива (колебания затухают);
        \item При $K = 5$ система находится на границе устойчивости (наблюдаются незатухающие колебания);
        \item При $K = 7$ система неустойчива (амплитуда нарастает по экспоненте).
    \end{itemize}
\end{enumerate}


\newpage
\section{СИНТЕЗ КОМАНДНОГО ГЕНЕРАТОРА}

\subsection{Аналитические расчеты}

Требуется спроектировать автономный командный генератор (КГ), способный воспроизводить заданный желаемый сигнал:
\[
    g_{\text{ж}}(t) = 9e^{-0.2t}\sin(3t + 5) + \sin(t).
\]

Рассмотрим составляющие желаемого выхода:
\begin{itemize}
    \item Гармоническая составляющая $\sin(t)$ с частотой $\omega_1 = 1$ рад/с соответствует паре чисто мнимых собственных значений:
    \[
        \lambda_{1,2} = \pm i.
    \]
    \item Затухающая гармоническая составляющая $9e^{-0.2t}\sin(3t + 5)$ с коэффициентом затухания $\sigma = 0.2$ и частотой $\omega_2 = 3$ рад/с соответствует паре комплексно-сопряженных собственных значений:
    \[
        \lambda_{3,4} = -0.2 \pm 3i.
    \]
\end{itemize}

Таким образом, матрица состояния $A$ должна иметь следующие собственные значения:
\[
    \sigma(A) = \{\pm i, \, -0.2 \pm 3i\}.
\]

Построим матрицу $A$ в блочно-диагональной форме (жордановой форме):
\[
    A = \begin{bmatrix}
        0 & 1 & 0 & 0 \\
        -1 & 0 & 0 & 0 \\
        0 & 0 & -0.2 & 3 \\
        0 & 0 & -3 & -0.2
    \end{bmatrix}.
\]

Общее решение однородной линейной системы при нулевом входе описывается выражением:
\[
    x(t) = e^{At}x(0), \qquad g(t) = C x(t) = C e^{At} x(0).
\]

Матричная экспонента $e^{At}$ для данной структуры имеет вид:
\[
    e^{At} = \begin{bmatrix}
        \cos(t) & \sin(t) & 0 & 0 \\
        -\sin(t) & \cos(t) & 0 & 0 \\
        0 & 0 & e^{-0.2t}\cos(3t) & e^{-0.2t}\sin(3t) \\
        0 & 0 & -e^{-0.2t}\sin(3t) & e^{-0.2t}\cos(3t)
    \end{bmatrix}.
\]

Выход системы $g(t)$ формируется как линейная комбинация элементов состояния:
\[
    g(t) = \begin{bmatrix} c_1 & c_2 & c_3 & c_4 \end{bmatrix} e^{At} x(0).
\]

Используя формулы приведения для тригонометрической функции желаемого сигнала с фазой $\sin(3t + 5) = \sin(3t)\cos(5) + \cos(3t)\sin(5)$, выберем параметры векторов так, чтобы в точности совпасть с заданным $g_{\text{ж}}(t)$. В результате получаем следующие параметры математической модели:
\begin{itemize}
    \item Матрица системы $A$:
    \[
        A = \begin{bmatrix}
            0 & 1 & 0 & 0 \\
            -1 & 0 & 0 & 0 \\
            0 & 0 & -0.2 & 3 \\
            0 & 0 & -3 & -0.2
        \end{bmatrix}
    \]
    \item Вектор начальных условий $x(0)$:
    \[
        x(0) = \begin{bmatrix}
            0 \\ 1 \\ 9\sin(5) \\ 9\cos(5)
        \end{bmatrix}
    \]
    \item Матрица выхода $C$:
    \[
        C = \begin{bmatrix} 1 & 0 & 1 & 0 \end{bmatrix}
    \]
\end{itemize}

\subsection{Моделирование работы командного генератора}

В соответствии с разработанной математической моделью была построена схема моделирования командного генератора в среде Simulink, использующая исключительно блоки элементарных операций (интеграторы, усилители и сумматоры).

\begin{figure}[H]
    \centering
    \includegraphics[width=0.95\textwidth]{scheme3.png}
    \caption{Схема моделирования командного генератора в среде Simulink}
    \label{fig:scheme3}
\end{figure}

На рисунках ниже представлены результаты симуляции работы командного генератора для Варианта 3: сравнительный график желаемого сигнала $g_{\text{ж}}(t)$ и выхода генератора $g(t)$, а также график погрешности $e(t) = g_{\text{ж}}(t) - g(t)$.

\begin{figure}[H]
    \centering
    \includegraphics[width=0.95\textwidth]{kg_comparison.png}
    \caption{Сравнение желаемого сигнала $g_{\text{ж}}(t)$ и выхода КГ $g(t)$}
    \label{fig:kg_comparison}
\end{figure}

\begin{figure}[H]
    \centering
    \includegraphics[width=0.95\textwidth]{kg_error.png}
    \caption{График ошибки моделирования $e(t) = g_{\text{ж}}(t) - g(t)$}
    \label{fig:kg_error}
\end{figure}

Как видно из полученных графиков, выходной сигнал автономного командного генератора с высокой точностью воспроизводит заданный желаемый сигнал. Ошибка моделирования находится в диапазоне машинной точности ($10^{-11}$), что подтверждает правильность аналитического синтеза и корректность построения схемы в Simulink.


\subsection{Выводы по заданию 3}
По результатам выполнения задания 3 можно сделать следующие выводы:
\begin{itemize}
    \item Успешно выполнена аналитическая разработка математической модели командного генератора в пространстве состояний для заданного сигнала $g_{\text{ж}}(t)$.
    \item Построена корректная схема моделирования в Simulink на основе элементарных операционных блоков.
    \item Компьютерное моделирование подтвердило высокую точность синтезированной системы: выходной сигнал $g(t)$ в точности воспроизводит желаемый сигнал с погрешностью в пределах машинной точности ($\sim 10^{-11}$).
\end{itemize}


\newpage
\section*{ЗАКЛЮЧЕНИЕ}
\addcontentsline{toc}{section}{ЗАКЛЮЧЕНИЕ}

В ходе выполнения лабораторной работы были исследованы основные динамические свойства линейных систем автоматического управления:
\begin{enumerate}
    \item Исследовано свободное движение линейных систем второго порядка при различных типах корней характеристического уравнения, подтверждена связь между их расположением и устойчивостью системы.
    \item Выполнен аналитический расчет и построена область устойчивости для системы третьего порядка с использованием критерия Гурвица, а также проверено влияние коэффициента передачи регулятора на характер переходных процессов.
    \item Осуществлен успешный аналитический синтез и компьютерное моделирование автономного командного генератора в среде Simulink, подтвердивший высокую точность воспроизведения заданного сложного гармонического сигнала.
\end{enumerate}
Все поставленные задачи выполнены в полном объеме, а результаты аналитических расчетов полностью подтвердились данными компьютерного моделирования.

\end{document}